% !TEX program = pdflatex
% THEME 1 — Difficile — Exercices
\documentclass[11pt,a4paper]{article}
\usepackage[utf8]{inputenc}
\usepackage[T1]{fontenc}
\usepackage{lmodern}
\usepackage[french]{babel}
\usepackage{amsmath,amssymb,mathtools}
\usepackage{icomma}
\usepackage{siunitx}
\sisetup{output-decimal-marker={,},per-mode=symbol}
\DeclareSIUnit{\atm}{atm}
\DeclareSIUnit{\bar}{bar}
\DeclareSIUnit{\mmHg}{mmHg}
\usepackage[version=4]{mhchem}
\usepackage{xcolor}
\usepackage{colortbl}
\usepackage{tikz}
\usepackage[most]{tcolorbox}
\usepackage{enumitem}
\usepackage{array}
\usepackage{booktabs}
\usepackage{fancyhdr}
\usepackage[a4paper,margin=2.1cm,headheight=26pt,headsep=10pt]{geometry}
\usetikzlibrary{arrows.meta,calc,shapes.geometric,positioning}

\definecolor{primary}{RGB}{150,98,162}
\definecolor{primarydark}{RGB}{92,52,110}
\definecolor{gazcol}{RGB}{0,120,190}
\definecolor{pcol}{RGB}{200,40,55}
\definecolor{tcol}{RGB}{198,93,9}

\newtcolorbox{enonce}[1][]{enhanced,breakable,colback=primary!4,colframe=primary,
  boxrule=0.8pt,arc=2pt,left=3mm,right=3mm,top=2mm,bottom=2mm,
  fonttitle=\bfseries,coltitle=white,title={#1}}

\pagestyle{fancy}\fancyhf{}
\fancyhead[L]{\small\textcolor{primarydark}{\textbf{Fiche 7} — Loi des gaz}}
\fancyhead[R]{\small\textcolor{primarydark}{\textsc{Thème 1 · Difficile · Exercices}}}
\fancyfoot[C]{\small\thepage}
\renewcommand{\headrulewidth}{0.8pt}
\renewcommand{\headrule}{\hbox to\headwidth{\color{primary}\leaders\hrule height \headrulewidth\hfill}}
\setlength{\parindent}{0pt}\setlength{\parskip}{3pt}

\newcounter{exo}
\newenvironment{exo}[1][]{\refstepcounter{exo}\medskip
  \noindent{\color{primarydark}\bfseries Exercice \theexo.}\ \textit{#1}\par\nopagebreak\smallskip}{\par}

\begin{document}

\begin{center}
{\Large\bfseries\color{primarydark} Thème 1 — Niveau Difficile}\\[2pt]
{\color{primary}\itshape Enchaîner plusieurs transformations et démasquer les pièges}
\end{center}
\textcolor{primary}{\hrule height 1pt}
\medskip

\begin{exo}[Un gaz suivi étape par étape dans un cylindre à piston]
Un cylindre muni d'un piston mobile contient une quantité fixée de gaz dans l'\textbf{état initial} :
$V_1=\qty{24}{\litre}$, $p_1=\qty{1.0}{\atm}$, $T_1=\qty{300}{\kelvin}$. On lui fait subir
successivement les deux étapes représentées ci-dessous.

\begin{center}
\begin{tikzpicture}[scale=0.9,>=Stealth,font=\footnotesize,
   stcap/.style={align=center,font=\scriptsize,anchor=north}]
  % Etat 1
  \draw[thick] (0,0) rectangle (1.6,3.0);
  \fill[gazcol!12] (0,0) rectangle (1.6,3.0);
  \draw[fill=primary!30,thick] (0,3.0) rectangle (1.6,3.25);
  \node[stcap] at (0.8,-0.35){\textbf{État 1}\\ $\qty{24}{\litre}$, $\qty{1}{\atm}$\\ $\qty{300}{\kelvin}$};
  % fleche etape A
  \draw[->,pcol,very thick] (2.0,1.6)--(3.4,1.6);
  \node[pcol] at (2.7,2.2){\textbf{Étape A}};
  \node[align=center] at (2.7,1.05){\scriptsize isotherme\\[-1pt]\scriptsize jusqu'à $\qty{2}{\atm}$};
  % Etat 2 (comprimé)
  \draw[thick] (3.8,0) rectangle (5.4,1.5);
  \fill[gazcol!12] (3.8,0) rectangle (5.4,1.5);
  \draw[fill=primary!30,thick] (3.8,1.5) rectangle (5.4,1.75);
  \node[stcap] at (4.6,-0.35){\textbf{État 2}\\ $V_A=?$, $\qty{2}{\atm}$\\ $\qty{300}{\kelvin}$};
  % fleche etape B
  \draw[->,tcol,very thick] (6.0,1.6)--(7.4,1.6);
  \node[tcol] at (6.7,2.2){\textbf{Étape B}};
  \node[align=center] at (6.7,1.05){\scriptsize isobare\\[-1pt]\scriptsize $300\to\qty{400}{\kelvin}$};
  % Etat 3 (chauffé, dilaté)
  \draw[thick] (7.8,0) rectangle (9.4,2.0);
  \fill[gazcol!12] (7.8,0) rectangle (9.4,2.0);
  \draw[fill=primary!30,thick] (7.8,2.0) rectangle (9.4,2.25);
  \node[stcap] at (8.6,-0.35){\textbf{État 3}\\ $V_B=?$, $\qty{2}{\atm}$\\ $\qty{400}{\kelvin}$};
\end{tikzpicture}
\end{center}
\vspace{0.3cm}

\begin{enumerate}[label=(\alph*),leftmargin=2em,itemsep=2pt]
  \item \textbf{Étape A} (isotherme) : le gaz est comprimé jusqu'à $p=\qty{2.0}{\atm}$ à température
        constante. Quelle loi appliques-tu ? Calcule le volume $V_A$ de l'état 2.
  \item \textbf{Étape B} (isobare) : à partir de l'état 2, on chauffe le gaz de $\qty{300}{\kelvin}$ à
        $\qty{400}{\kelvin}$, la pression restant à $\qty{2.0}{\atm}$. Calcule le volume $V_B$ de
        l'état 3.
  \item Retrouve directement $V_B$ en appliquant la \textbf{loi combinée} entre l'état 1 et l'état 3,
        sans passer par l'état 2. Le résultat est-il cohérent avec (b) ?
  \item La quantité $\dfrac{pV}{T}$ est-elle la même dans les trois états ? Que représente-t-elle ?
\end{enumerate}
\end{exo}

\begin{exo}[Une mesure en unités du Système international]
Lors d'une expérience, un gaz exerce une pression $p_1=\qty{9.6e5}{\pascal}$ alors qu'il occupe un
volume $V_1=\qty{0.035}{\meter\cubed}$ à la température $T_1=\qty{220}{\kelvin}$. On réduit son volume à
$V_2=\qty{0.0035}{\meter\cubed}$ et on porte sa température à $T_2=\qty{440}{\kelvin}$
(la quantité de gaz ne change pas).
\begin{enumerate}[label=(\alph*),leftmargin=2em,itemsep=2pt]
  \item Parmi $p$, $V$, $T$ et $n$, lesquelles varient au cours de cette transformation ? En déduis la
        loi à utiliser.
  \item Établis l'expression littérale de la pression finale $p_2$.
  \item Calcule $p_2$ (en pascals). Toutes les grandeurs étant en unités SI, aucune conversion n'est
        nécessaire : explique pourquoi.
  \item La pression a-t-elle augmenté ou diminué, et de quel facteur ?
\end{enumerate}
\end{exo}

\begin{exo}[Doubler une pression\dots\ et retrouver le zéro absolu]
Un récipient \textbf{rigide} contient un gaz sous $p_1=\qty{1.0}{\atm}$ à $\theta_1=\qty{0}{\degreeCelsius}$.
\begin{enumerate}[label=(\alph*),leftmargin=2em,itemsep=2pt]
  \item À quelle température (en \si{\degreeCelsius}) faut-il porter le gaz pour \textbf{doubler} sa
        pression, à volume constant ?
  \item Un étudiant propose : « pour doubler la pression, il suffit de doubler la température, soit
        $2\times\qty{0}{\degreeCelsius}=\qty{0}{\degreeCelsius}$, ce qui est impossible ». Où est son
        erreur de raisonnement ?
  \item En prolongeant la logique de la loi isochore, à quelle température (en \si{\degreeCelsius}) la
        pression du gaz \emph{s'annulerait-elle} théoriquement ? Quel nom porte cette température
        remarquable ?
\end{enumerate}
\end{exo}

\end{document}
